%!TEX root = ../template.tex
%%%%%%%%%%%%%%%%%%%%%%%%%%%%%%%%%%%%%%%%%%%%%%%%%%%%%%%%%%%%%%%%%%%%
%% abstract-pt.tex
%% NOVA thesis document file
%%
%% Abstract in Portuguese
%%%%%%%%%%%%%%%%%%%%%%%%%%%%%%%%%%%%%%%%%%%%%%%%%%%%%%%%%%%%%%%%%%%%

\typeout{NT FILE abstract-pt.tex}%

A modelação de software e a modelação conceptual são frequentemente incluídas em cursos de Engenharia de Software e áreas relacionadas, mas tendem a receber menos destaque do que outros tópicos na educação em Computação. Os estudantes frequentemente percecionam a modelação como um apoio essencial à compreensão e comunicação ou, pelo contrário, como documentação extra. A investigação existente fornece contributos relevantes, mas encontra-se distribuída por diferentes tradições de modelação, com explicação integrada ainda limitada sobre o modo como os desafios emergem e se relacionam com estratégias, apoio e condições da unidade curricular.

Esta tese investiga como a modelação de software é experienciada, ensinada e aprendida no ensino superior, tratando a modelação como uma atividade socio-técnica. O estudo é guiado por três questões: (RQ1) que desafios estudantes e docentes percecionam, e como se manifestam ao interpretar, criar, refinar e avaliar modelos; (RQ2) que estratégias e necessidades de apoio (pedagógicas e de ferramentas) são identificadas; e (RQ3) de que forma ferramentas, desenho curricular, domínios e práticas de feedback moldam essas experiências.

O estudo usou Socio-Technical Grounded Theory. Catorze entrevistas semi-estruturadas (quatro docentes, dez estudantes) foram sujeitas a codificação aberta (277 códigos e 78 memorandos) e comparadas em 14 categorias e quatro fenómenos. A amostragem teórica da segunda vaga densificou propriedades sem acrescentar uma décima quinta categoria de visão geral. A análise argumenta saturação local desse conjunto de categorias e suficiência teórica para esta amostra, não saturação teórica global.

O relato mostra um primeiro passo de modelação condicionado por um caminho trabalhado, distinto da sobrecarga de símbolos; dificuldade em julgar ``suficientemente bom'' sem um resultado de compilação; feedback tardio e de acesso desigual, com retrabalho dependente do tempo restante e da nota já fechada; procura de verificações e explicações em vez de trabalhos gerados automaticamente; e diagramas que só se mantêm quando alguém os assume e alguém os lê. A carga curricular e o atrito das ferramentas condicionam a existência efetiva de exemplos e ciclos de feedback. As implicações para o ensino e para ferramentas educativas derivam destes padrões e não foram ensaiadas.

% Palavras-chave do resumo
\keywords{
  educação em modelação de software \and
  modelação conceptual \and
  grounded theory \and
  perspetiva socio-técnica \and
  competência de modelação
}